\section{EXPERIMENTAL RESULTS}
\subsection{Experimental Design}

\begin{frame}[t]{Experimental Design}
\scriptsize
\vspace{-0.12cm}
\begin{block}{Evaluation Protocol}
Models were evaluated using repeated train--test splits across five random seeds: 42, 52, 62, 72, and 82.
\end{block}

\vspace{0.08cm}
\begin{table}
\centering
\renewcommand{\arraystretch}{1.04}
\setlength{\tabcolsep}{3pt}
\begin{tabularx}{\textwidth}{p{3.0cm} p{2.4cm} p{2.0cm} p{2.0cm} X}
\toprule
\textbf{Dataset} & \textbf{Task} & \textbf{Original Size} & \textbf{Rows Used} & \textbf{Features} \\
\midrule
Heart Disease & Classification & 1,025 & 302 & 13 \\
Diabetes & Classification & 768 & 768 & 8 \\
Loan Prediction & Classification & 614 & 614 & 11 \\
Cyber Security Attacks & Classification & 40,000 & 5,000 & 18 \\
Abalone & Regression & 4,177 & 4,177 & 8 \\
California Housing & Regression & 20,640 & 5,000 & 12 \\
Wind Turbine & Regression & 50,530 & 5,000 & 3 \\
\bottomrule
\end{tabularx}
\end{table}

\vspace{0.05cm}
\begin{columns}[T,totalwidth=\textwidth]
\begin{column}{0.48\textwidth}
\textbf{Classification:} Accuracy, Precision, Recall, F1-score.
\end{column}
\begin{column}{0.48\textwidth}
\textbf{Regression:} MAPE, RMSE, $R^2$, average target value.
\end{column}
\end{columns}
\end{frame}

\subsection{Classification Results}
\begin{frame}[t]{Classification Results}
\scriptsize
\begin{table}
\centering
\renewcommand{\arraystretch}{1.16}
\setlength{\tabcolsep}{3.8pt}
\resizebox{\textwidth}{!}{%
\begin{tabular}{lccccc}
\toprule
\textbf{Dataset} & \textbf{KFCCC} & \textbf{LR} & \textbf{RFC} & \textbf{SVM} &  \textbf{CC} \\
\midrule
Heart Disease & \cellcolor{green!20}\textbf{0.819 $\pm$ 0.037} & 0.806 $\pm$ 0.032 & 0.806 $\pm$ 0.060 & 0.800 $\pm$ 0.042 & 0.768 $\pm$ 0.062 \\
Diabetes & 0.757 $\pm$ 0.056 & 0.755 $\pm$ 0.047 & \cellcolor{green!20}\textbf{0.761 $\pm$ 0.036} & 0.753 $\pm$ 0.047 & 0.751 $\pm$ 0.050 \\
Loan Prediction & 0.820 $\pm$ 0.030 & \cellcolor{green!20}\textbf{0.826 $\pm$ 0.024} & 0.795 $\pm$ 0.023 & 0.824 $\pm$ 0.024 & 0.821 $\pm$ 0.028 \\
Cyber Security & 0.339 $\pm$ 0.013 & \cellcolor{green!20}\textbf{0.340 $\pm$ 0.006} & 0.333 $\pm$ 0.012 & \cellcolor{green!20}\textbf{0.340 $\pm$ 0.007} & 0.334 $\pm$ 0.004 \\
\bottomrule
\end{tabular}}
\end{table}

\vspace{0.12cm}
\begin{block}{Key Interpretation}
\begin{itemize}\setlength\itemsep{0.06cm}
    \item \textbf{KFCCC} performs best on Heart Disease and remains competitive overall.
    \item Baselines are stronger on some datasets, showing that clusterwise learning depends on dataset structure.
    \item The result supports the need to compare KFC, COBRA, and standard baselines in one package.
\end{itemize}
\end{block}

{\scriptsize \textbf{Note:} KFCCC = KFCClassifier-CombinedClassifier; LR = Logistic Regression; RFC = Random Forest Classifier; SVM = Support Vector Machine; CC = CombinedClassifier. Green cells indicate best mean accuracy.}
\end{frame}

\subsection{Regression Results}
\begin{frame}[t]{Regression Results}
\scriptsize
\begin{table}
\centering
\renewcommand{\arraystretch}{1.22}
\setlength{\tabcolsep}{5pt}
\resizebox{\textwidth}{!}{%
\begin{tabular}{lccccc}
\toprule
\textbf{Dataset} & \textbf{KFCGC} & \textbf{MCR} & \textbf{GC} & \textbf{RFR} & \textbf{GBR} \\
\midrule
Abalone & \cellcolor{green!20}\textbf{0.560 $\pm$ 0.012} & 0.555 $\pm$ 0.018 & 0.550 $\pm$ 0.019 & 0.547 $\pm$ 0.021 & 0.537 $\pm$ 0.020 \\
California Housing & 0.778 $\pm$ 0.009 & 0.805 $\pm$ 0.064 & 0.791 $\pm$ 0.054 & \cellcolor{green!20}\textbf{0.874 $\pm$ 0.007} & 0.858 $\pm$ 0.008 \\
Wind Turbine & 0.843 $\pm$ 0.013 & \cellcolor{green!20}\textbf{0.876 $\pm$ 0.009} & 0.874 $\pm$ 0.009 & 0.872 $\pm$ 0.005 & \cellcolor{green!20}\textbf{0.876 $\pm$ 0.008} \\
\bottomrule
\end{tabular}}
\end{table}

\vspace{0.12cm}
\begin{block}{Key Interpretation}
\begin{itemize}\setlength\itemsep{0.06cm}
    \item \textbf{KFCGC} is strongest on Abalone and competitive on other regression datasets.
    \item Random Forest and gradient-based baselines remain stronger on some larger datasets.
    \item KFCGC is promising, but its benefit depends on meaningful subgroup structure.
\end{itemize}
\end{block}

{\scriptsize \textbf{Note:} KFCGC = KFCRegressor-GradientCOBRA; MCR = MixCOBRARegressor; GC = GradientCOBRA; RFR = Random Forest Regressor; GBR = Gradient Boosting Regressor. Green cells indicate best mean $R^2$.}
\end{frame}
